\section{Results}\label{sec:results}
All reported results are derived from deterministic REPLAY of recorded evaluator outputs unless explicitly noted. All sessions within a condition are dependent across question orders due to the 8-cluster structure (Sect.~\ref{sec:method}).

\noindent\textbf{Trajectory Divergence and Reduced Level Changes (RQ1).} Enabling the agent layer ($M_1$) altered the difficulty trajectory relative to the frozen baseline ($B_0$) in \textbf{all 160 replay sessions (100\%)} (A7; Table~\ref{tab:xapaired}). The $M_1$ trajectories show fewer level changes and reversals than $B_0$: mean level changes fell from 15.12 to 12.45 per session (point difference $-2.67$, naive session bootstrap 95\% CI $[-3.19, -2.13]$), and mean reversals fell from 11.74 to 9.87 per session ($-1.88$, naive CI $[-2.28, -1.46]$). Agreement with the uncoordinated PPO action dropped from 0.626 to 0.400 (point difference $-0.226$, cluster CI \phagrci{}), with Strategy proposals winning pre-guardrail arbitration in 11.9 decisions per session and remaining the final unoverridden action source in 6.669 decisions after downstream guardrail evaluation (\textbf{strat.} in Table~\ref{tab:xa}).


\begin{table}[tbp]
\centering\scriptsize\setlength{\tabcolsep}{3pt}
\caption{Replay ablation summary across 160 sessions per condition (mean per session unless noted as counts). \textbf{compl./auth./trans.}: sessions completing, keeping committed scores authoritative, and preserving valid transitions. \textbf{chg./rev.}: mean level changes and reversals; \textbf{agr.}: fraction of decisions agreeing with the PPO action; \textbf{gfire}: guardrail overrides; \textbf{strat.}: decisions where the Strategy Agent was the final action source; \textbf{byp.}: sessions with guardrail bypass; \textbf{p50/p95}: pooled per-turn latency in ms; \textbf{calls}: agent calls per session.}\label{tab:xa}
\resizebox{\textwidth}{!}{\begin{tabular}{lcccrrrrrrrrr}
\toprule
Cond. & compl. & auth. & trans. & chg. & rev. & agr. & gfire & strat. & byp. & p50 & p95 & calls\\
\midrule
B0 & 160 & 160 & 160 & 15.119 & 11.744 & 0.626 & 9.588 & 0.0 & 0 & 0.546 & 2.056 & -- \\
B7 & 160 & 160 & 160 & 2.675 & 1.256 & -- & 3.831 & 0.0 & 0 & 0.463 & 1.059 & -- \\
B0c & 160 & 160 & 160 & 15.119 & 11.744 & 0.626 & 9.588 & 0.0 & 0 & 0.553 & 2.004 & -- \\
M1 & 160 & 160 & 160 & 12.45 & 9.869 & 0.4 & 6.838 & 6.669 & 0 & 1.589 & 4.411 & 47.68 \\
M2 & 160 & 160 & 160 & 15.119 & 11.744 & 0.626 & 9.588 & 0.0 & 0 & 1.218 & 3.513 & 23.25 \\
M3 & 160 & 160 & 160 & 12.45 & 9.869 & 0.4 & 6.838 & 6.669 & 0 & 1.118 & 3.189 & 23.84 \\
M4 & 160 & 160 & 160 & 2.225 & 1.0 & 0.014 & 0.0 & 19.619 & 3 & 1.573 & 4.357 & 47.68 \\
M5 & 160 & 160 & 160 & 12.45 & 9.869 & 0.4 & 6.838 & 6.669 & 0 & 1.581 & 4.553 & 47.68 \\
M6 & 160 & 160 & 160 & 12.45 & 9.869 & 0.4 & 6.838 & 6.669 & 0 & 1.592 & 4.46 & 47.68 \\
M8 & 160 & 160 & 160 & 12.45 & 9.869 & 0.4 & 6.838 & 6.669 & 0 & 1.575 & 4.35 & 47.68 \\
M7 & 160 & 160 & 160 & 2.706 & 1.312 & -- & 3.812 & 9.037 & 0 & 1.502 & 2.968 & 47.68 \\
M1c & 160 & 160 & 160 & 12.45 & 9.869 & 0.4 & 6.838 & 6.669 & 0 & 1.573 & 4.534 & 47.68 \\
\bottomrule
\end{tabular}
}
\end{table}

\begin{table}[tbp]
\centering\scriptsize\setlength{\tabcolsep}{3pt}
\caption{Paired differences ($A - B$) over identical 160 replay sessions: fraction of sessions with divergent difficulty paths, paired mean differences with naive session-level 95\% percentile-bootstrap intervals, and pooled median turn latency delta. Intervals are descriptive due to cluster reuse.}\label{tab:xapaired}
\resizebox{\textwidth}{!}{\begin{tabular}{lrcccc}
\toprule
A$-$B & path$\ne$ & $\Delta$ level chg. & $\Delta$ reversals & $\Delta$ PPO agree & $\Delta$ p50 ms\\
\midrule
M1$-$B0 & 1.0 & $-$2.669 [$-$3.188,$-$2.131] & $-$1.875 [$-$2.275,$-$1.462] & $-$0.226 [$-$0.244,$-$0.208] & +1.043 [+1.000,+1.079] \\
M2$-$B0 & 0.0 & 0 & 0 & 0 & +0.672 [+0.646,+0.698] \\
M3$-$M1 & 0.0 & 0 & 0 & 0 & $-$0.471 [$-$0.511,$-$0.425] \\
M4$-$M1 & 1.0 & $-$10.225 [$-$10.544,$-$9.900] & $-$8.869 [$-$9.275,$-$8.463] & $-$0.387 [$-$0.403,$-$0.370] & $-$0.016 [$-$0.065,+0.033] \\
M5$-$M1 & 0.0 & 0 & 0 & 0 & $-$0.008 [$-$0.059,+0.037] \\
M6$-$M1 & 0.0 & 0 & 0 & 0 & +0.004 [$-$0.030,+0.053] \\
M8$-$M1 & 0.0 & 0 & 0 & 0 & $-$0.013 [$-$0.057,+0.032] \\
M7$-$B7 & 0.025 & +0.031 [+0.000,+0.081] & +0.056 [+0.006,+0.125] &  & +1.039 [+0.996,+1.070] \\
B0c$-$B0 & 0.0 & 0 & 0 & 0 & +0.008 [$-$0.019,+0.027] \\
M1c$-$M1 & 0.0 & 0 & 0 & 0 & $-$0.016 [$-$0.055,+0.032] \\
\bottomrule
\end{tabular}
}
\end{table}

\noindent\textbf{Cluster Dependence, Fragility, and Proactive Probing Sensitivity.} Because the 160 sessions reuse 8 underlying replay clusters, treating sessions as independent underestimates parameter uncertainty. A post-hoc cluster bootstrap over the 8 replay clusters (Field and Welsh~\cite{fw07}, $B = 10{,}000$) yields wider, defensible intervals: level changes \phchgmean{} \phchgci{}, reversals \phrevmean{} \phrevci{}, and PPO action agreement \phagrmean{} \phagrci{}. The negative sign for level changes and reversals is shared across 7 of the 8 replay clusters.

Proactive probing added an average of \phfumean{} turns per session (naive CI $[+0.34, +0.85]$, cluster CI \phfuci{}). However, cluster-level decomposition reveals that this probe increase is present in \textbf{only \phfusame{} of the 8 replay clusters} ($k=1$); in the remaining seven clusters, proactive probing added exactly 0.00 turns. A leave-one-cluster-out (LOCO) sensitivity analysis confirms that omitting cluster $k=1$ shifts the level-change delta to $-3.76$ (CI $[-4.81, -2.78]$), whereas omitting any of the other seven clusters causes the 95\% interval to include zero. Thus, while a directional reduction in level changes and reversals is observed across most replay clusters (7 of 8) and the primary cluster-level interval excludes zero, the magnitude is sensitive to the small set of structural clusters, and extra-turn gains cannot be generalized as an unconditional property.

\noindent\textbf{Proposal Distribution.} Across the 160 sessions, the Strategy Agent generated \scTotal{} proposals: \scRemEight{} \emph{remediate} proposals (confidence 0.80), \scRecEight{} and \scRecSeven{} \emph{recover} proposals (confidence 0.80 and 0.70), and \scSteady{} \emph{steady} proposals (confidence 0.45). Notably, the Strategy Agent emitted zero \emph{stretch} proposals (\scStretch{}), indicating that during replay the candidate trajectories never maintained the sustained high-performance criteria ($3\text{ scores} \ge 0.60$, $\text{EMA} \ge 0.70$) required to trigger proactive difficulty increases.
\noindent\textbf{Coordination Overhead.} Introducing the multi-agent layer increased pooled per-turn latency from \xaBzeroMed{}\,ms ($B_0$) to \xaMoneMed{}\,ms ($M_1$) at the median (end-to-end orchestration overhead $+1.04$\,ms, 95\% CI $[+1.00, +1.08]$\,ms) and from \xaBzeroHigh{}\,ms to \xaMoneHigh{}\,ms at the 95th percentile, executing an average of 47.7 agent calls per session.


\begin{table}[tbp]
\centering\scriptsize\setlength{\tabcolsep}{3pt}
\caption{Registered predictions and empirical outcomes for Experiment XA.}\label{tab:preds}
\begin{tabular}{l>{\raggedright\arraybackslash}p{7.6cm}l}
\toprule
ID & Prediction & Outcome\\
\midrule
A1 & completion 100\% in every condition & Supported \\
A2 & committed scores equal the recorded evaluator output in every session of every condition & Supported \\
A3 & difficulty levels stay within [1,5], and every RL-phase decision moves by at most one level, in every session of every condition & Supported \\
A4 & M2 (no Strategy Agent) reproduces B0 (difficulty, score and question paths) in every session & Supported \\
A5 & M3, M5, M6, M8 reproduce M1's difficulty/score/question paths in every session & Supported \\
A6 & M4 (no arbitration) has guardrail\_bypass>0 in $\ge$1 session; M1 has 0 in all & Supported \\
A7 & M1 differs from B0 in difficulty path in $\ge$50\% of sessions & Supported \\
A8 & median per-turn latency M1 $-$ B0 < 50 ms (pooled turns) & Supported \\
A9 & with PPO unavailable, M7 differs from B7 in difficulty path in $\ge$50\% of sessions & Contradicted \\
A10 & canonical guardrails change the difficulty path in $\ge$1 session vs legacy (B0c vs B0) & Contradicted \\
\bottomrule
\end{tabular}

\end{table}

\noindent\textbf{Decomposing Arbitration via Threshold Sensitivity (RQ3).} Disabling the Strategy Agent ($M_2$) reproduced the baseline ($B_0$) difficulty, score, and question trajectories identically in 160 of 160 sessions (A4), so the coordination layer leaves trajectories unchanged when Strategy is disabled and Coaching remains active. In the unarbitrated condition ($M_4$), where all arbitration checks were simultaneously disabled, trajectory dynamics collapsed: level changes fell to 2.225 per session, PPO agreement decreased to 0.014, and guardrail bypass occurred (Table~\ref{tab:xa}). Experiment XC (Sect.~\ref{sec:method}) isolates the specific gating mechanism by varying only the confidence threshold $\tau$ from 0.0 to 0.9 while keeping the other validation, evidence, and guardrail conditions active (Table~\ref{tab:xc}).


\begin{table}[tbp]
\centering\scriptsize\setlength{\tabcolsep}{3pt}
\caption{Experiment XC (REPLAY, 160 sessions per row, mean per session): only confidence threshold $\tau$ varies; evidence requirements and guardrails remain active. \textbf{chg./rev.}: mean level changes and reversals; \textbf{agr.}: agreement with PPO; \textbf{strat.}: decisions from Strategy Agent; \textbf{gfire}: guardrail overrides; \textbf{byp.}: sessions with guardrail bypass.}\label{tab:xc}
\begin{tabular}{rcccccc}\toprule
$\tau$ & chg. & rev. & agr. & strat. & gfire & byp.\\
\midrule
0.0 & 2.225 & 1.000 & 0.014 & 16.075 & 3.544 & 0 \\
0.5 & 12.450 & 9.869 & 0.400 & 6.669 & 6.838 & 0 \\
0.6 & 12.450 & 9.869 & 0.400 & 6.669 & 6.838 & 0 \\
0.7 & 12.450 & 9.869 & 0.400 & 6.669 & 6.838 & 0 \\
0.8 & 12.450 & 9.819 & 0.404 & 6.594 & 6.863 & 0 \\
0.9 & 15.119 & 11.744 & 0.626 & 0.000 & 9.588 & 0 \\
\bottomrule
\end{tabular}

\end{table}

At $\tau = 0.0$, the difficulty trajectory collapses toward the arbitration-ablated regime (\xctauzerochg{} level changes, 1.000 reversals, and \xctauzeroagr{} PPO agreement) while guardrail bypass remains zero in all 160 sessions. Because unevidenced \emph{steady} proposals are assigned confidence 0.45, setting $\tau = 0.0$ allows uncalibrated default proposals to override PPO on most RL decisions (16.1 of 19.6 per session), indicating that proposal admission through confidence-gated arbitration is a major control point for the observed trajectory under the tested replay conditions. Behavior remained flat across $\tau \in [0.5, 0.7]$ (12.450 level changes, 0.400 agreement; 0.404 agreement at $\tau = 0.8$) because realized evidenced confidences occur at discrete values (0.70 and 0.80). At $\tau = 0.9$, the threshold exceeds all realized heuristic confidences ($\max = 0.80$), preventing Strategy proposals from superseding PPO and reverting trajectories to the baseline ($B_0$). All \xcsupported{} registered XC predictions were supported.

\noindent\textbf{Peer Information Exchange and Grounding Ablations ($M_5, M_6$).} Ablating grounding verification ($M_5$) or the unidirectional peer message from Strategy to Coaching ($M_6$) resulted in difficulty trajectories identical to the full configuration ($M_1$) across all 160 sessions (A5). The peer message does not feed back into Strategy decision-making; its role is strictly to condition Coaching presentation (switching depth to \emph{detailed} during \emph{recover}, \emph{brief} during \emph{stretch}, and attaching remediation tags during \emph{remediate}).

\noindent\textbf{Contradicted Prediction A9 (PPO Fallback Interaction).} Registered prediction A9 posited that disabling the learned PPO policy would alter $\ge 50\%$ of session paths. Empirically, disabling PPO ($M_7$) altered difficulty trajectories in only \textbf{4 of 160 sessions (2.5\%)}. Inspection of orchestrator traces reveals that under replay conditions, the rule-based Strategy Agent heuristic predominantly reaches the identical difficulty action as the non-RL heuristic fallback.

\noindent\textbf{Contradicted Prediction A10 (Guardrail Divergence).} Registered prediction A10 posited that evaluating the system with canonical guardrails ($B_{0c}$) would alter at least one trajectory path. In replay, \textbf{zero paths differed (0/160)}. While the two implementations disagree on 11.8\% of synthetic state grid points (4,728 of 39,930), this disagreement is a property of the synthetic grid rather than an observed replay or deployment frequency. The divergent grid regions correspond to extreme acoustic hesitation or low voice-confidence states; these states were not reached in the evaluated replay trajectories, where acoustic defaults remain neutral.

\noindent\textbf{Trajectory Invariance under Functional Decomposition (RQ2).} A central motivation for multi-agent architecture is the assumption that role specialization produces different decision behavior. To illustrate this assumption, we examine Condition $M_3$, whose outcome is largely fixed by the architecture. The implementation exhibits a structural property implied by the architecture's authority and information topology: \texttt{CoachingPlan} contains no difficulty-control fields, \texttt{arbitrate\_difficulty()} does not inspect the coaching plan, and inter-agent communication is unidirectional from Strategy to Coaching through \texttt{TurnObservation.peer}. Under this architecture, Coaching execution does not provide an information pathway back to the Strategy difficulty decision.

Consistent with this structural property, the matched replay experiment produced identical difficulty trajectories in all 160 evaluated sessions ($M_3 \equiv M_1$, Table~\ref{tab:xapaired}), with identical level changes (12.45), directional reversals (9.87), and score progressions (A5).

\begin{table}[tbp]
\centering\scriptsize\setlength{\tabcolsep}{3pt}
\caption{Architectural and operational comparison of the Specialized Two-Agent Architecture ($M_{\text{split}}$) versus a Unified Monolithic Controller ($M_{\text{unified}}$).}\label{tab:specialization}
\begin{tabular}{lp{3.7cm}p{3.7cm}}
\toprule
\textbf{Dimension} & \textbf{Specialized Architecture ($M_{\text{split}}$)} & \textbf{Unified Architecture ($M_{\text{unified}}$)} \\
\midrule
\textbf{Agent Roles} & Strategy Agent + Coaching Agent & Unified Technical-Coach Agent \\
\textbf{State Encapsulation} & Separate private ledgers (Weakness vs Coaching) & Monolithic shared session ledger \\
\textbf{Inter-Agent Comm.} & Explicit typed peer message (\texttt{TurnObservation.peer}) & Internal variable passing \\
\textbf{Decision Equivalence} & Baseline reference ($M_1$, 160 sessions) & $\mathbb{P} = 1.0$ ($1{,}280/1{,}280$ turns identical) \\
\textbf{Serialization Micro-bench.} & 2 models ($0.0092$\,ms per turn) & 1 model ($0.0038$\,ms, $-58.7\%$) \\
\textbf{Component Dispatch Latency} & $0.0440$\,ms per dispatch & $0.0401$\,ms ($-8.9\%$) \\
\textbf{End-to-End Orchestration} & $+1.04$\,ms p50 ($M_1$ vs $B_0$ only) & Not measured \\
\textbf{Decision-record structure} & Separate per-agent decision records & Joint session decision state \\
\bottomrule
\end{tabular}

\end{table}

\noindent\textbf{Decision Equivalence under Matched Replay.} On identical observation inputs under matched heuristic rules, a standalone unified controller (\texttt{UnifiedCoachAgent}, a joint-ledger re-implementation of the same rule set) synthesizing both strategy and coaching outputs produced identical normalized decision outputs across all 160 sessions (1,280 turns):
$$\mathbb{P}(\text{Decision}_{\text{unified}} = \text{Decision}_{\text{split}}) = 1.0$$
For this comparison, a turn-level decision denotes the normalized decision tuple: Strategy difficulty action ($\{\text{Easier}, \text{Same}, \text{Harder}\}$), mode, confidence, and probe concept, along with Coaching depth, hint/example flags, tone, remediation tag, and emphasis concepts. Under matched replay, architectural decomposition did not alter observed decision outputs in this evaluation setting.

\noindent\textbf{Architectural Characteristics and Operational Overhead.} As synthesized in Table~\ref{tab:specialization}, specialization under the tested configuration establishes separate state ledgers without altering the observed decision outputs, while introducing measured operational overhead:
\begin{itemize}\setlength{\itemsep}{1pt}\setlength{\parskip}{0pt}
\item \emph{State Decoupling (Architectural Design Property):} Strategy and Coaching maintain separate memory structures with separate schema boundaries; this separation yields no decision divergence in replay.
\item \emph{Operational Cost (Measured Empirical Outcome):} Runtime overhead is reported at three distinct measurement scopes; only (ii) and (iii) compare $M_{\text{split}}$ with $M_{\text{unified}}$: (i)~\emph{End-to-end orchestration overhead} ($M_1$ vs $B_0$; not measured for split vs unified): $+1.04$\,ms median turn latency ($B_0$ 0.55\,ms vs $M_1$ 1.59\,ms; $+2.36$\,ms p95; 47.7 calls/session); (ii)~\emph{Component dispatch latency}: a narrower dispatch measurement of $0.0440$\,ms per turn for split vs $0.0401$\,ms for unified ($-8.9\%$, saving $0.0039$\,ms; descriptive, single-run measurements); and (iii)~\emph{Serialization micro-benchmark}: validating two Pydantic models requires $0.0092$\,ms vs $0.0038$\,ms for one model ($-58.7\%$, saving $0.0054$\,ms).\end{itemize}

\noindent\textbf{Authority Boundaries, Invariant Enforcement, and Safety Containment (RQ4).} Every evaluated condition completed all 160 sessions, maintained strict evaluator authority (committed scores equaled recorded evaluator output in 100\% of sessions), and enforced valid difficulty bounds in $[1, 5]$ with transitions of at most one level per step (A1--A3, Table~\ref{tab:preds}; exact Clopper--Pearson 95\% lower bound for 160/160: 0.977~\cite{cp34}; treating sessions as independent; descriptive). In the unarbitrated condition ($M_4$), where checks were removed, guardrail bypass occurred; in contrast, under full bounded authority ($M_1$) and throughout the threshold decomposition (XC), guardrail bypass remained at zero in all 160 sessions. Under targeted fault injection (Experiment XB; 40 sessions per cell, faults fired at rate 0.3 per call), global score authority and legal transition bounds were preserved in 160/160 sessions across all variants. When adversarial-valid proposals opposing the evidence were injected, 38 of 40 sessions bypassed guardrails without arbitration (95\% Clopper--Pearson CI $[0.831, 0.994]$), whereas full arbitration contained bypass to 0 of 40 sessions (95\% CI $[0.000, 0.088]$). The per-agent circuit breaker opened in 29 of 40 sessions (72.5\%, 95\% CI $[0.561, 0.854]$), falling below the registered 80\% threshold. Timing-out agent proposals produced a pooled p95 turn latency of 154.4\,ms; the individual agent call was bounded by the configured 150\,ms timeout, and the turn latency additionally includes orchestration overhead. Ungrounded concepts were emphasized in 40 of 40 sessions when the grounding check was removed, but in 0 of 40 with it active. However, in the separate lexical probe (Experiment XD), unsupported paraphrased credit was accepted in 57 of 57 applicable authored cases (95\% CI $[0.937, 1.000]$), demonstrating that structural and lexical validation enforces an evidence contract rather than semantic verification. These results establish containment of the tested fault classes only.
